% chapters/04/contoh.tex
\section{Dua Contoh Terapan}
\label{sec:contoh-bab-4}

\subsection{Contoh Pertama: Anggaran Waktu Satu Bingkai}

\begin{contohbox}
Sebuah program dijalankan dengan batas laju 60 bingkai per detik. Jatah waktu
untuk satu bingkai adalah

\[
\frac{1000\ \text{milidetik}}{60} = 16{,}7\ \text{milidetik}.
\]

Misalkan pekerjaan program terbagi sebagai berikut: membaca masukan 0,5
milidetik, memperbarui keadaan 5 milidetik, dan menggambar 25 milidetik.
Jumlahnya 30,5 milidetik, jauh melebihi jatah 16,7 milidetik.

Karena jatahnya sudah terlampaui, \code{clock.tick()} tidak menahan apa pun.
Siklus berjalan secepat kemampuan program, yaitu satu bingkai setiap 30,5
milidetik, atau sekitar 32,8 bingkai per detik. Batas 60 yang diminta tidak
pernah tercapai.
\end{contohbox}

Tabel~\ref{tab:anggaran-bingkai} memperlihatkan akibatnya pada tiga nilai
batas laju, dengan pekerjaan yang sama beratnya.

\begin{table}[htbp]
\centering
\caption{Laju tercapai pada tiga batas laju dengan pekerjaan 30,5 milidetik per bingkai}
\label{tab:anggaran-bingkai}
{\setlength{\tabcolsep}{4pt}\footnotesize%
\begin{tabularx}{\textwidth}{@{}C{2.0cm} C{2.8cm} C{2.8cm} Y@{}}
\toprule
\headrow \textbf{Batas laju} & \textbf{Jatah per bingkai} & \textbf{Laju tercapai} & \textbf{Keterangan} \\
\midrule
30 & 33,3 milidetik & 30,0 FPS & Pekerjaan masih muat, sisa 2,8 milidetik ditahan. \\
\addlinespace
60 & 16,7 milidetik & 32,8 FPS & Jatah terlampaui; batas tidak dapat dipenuhi. \\
\addlinespace
120 & 8,3 milidetik & 32,8 FPS & Menaikkan batas tidak mengubah apa pun. \\
\bottomrule
\end{tabularx}}
\end{table}

Contoh ini memuat satu pelajaran yang penting: \textbf{menaikkan batas laju
tidak mempercepat program yang pekerjaannya memang berat}. Pada batas 60 dan
120, laju yang tercapai sama, yaitu 32,8 bingkai per detik, karena yang
menentukan adalah lama pekerjaannya, bukan batas yang diminta.

Apabila Anda menginginkan laju yang lebih tinggi, yang perlu dikurangi adalah
pekerjaan satu bingkai, bukan batasnya. Pada contoh di atas, tahap
penayangan memakan 25 dari 30,5 milidetik. Mengurangi pekerjaan penayangan,
misalnya dengan menggambar lebih sedikit objek atau memakai gambar yang lebih
kecil, memberi perbaikan yang jauh lebih besar daripada menaikkan batas.

\subsection{Contoh Kedua: Membaca Gejala dan Menentukan Penyebabnya}

Gejala yang tampak di layar jarang menyebutkan penyebabnya secara langsung.
Tabel~\ref{tab:gejala-siklus} memuat gejala yang paling sering muncul pada
siklus permainan beserta cara memeriksanya.

\begin{table}[htbp]
\centering
\caption{Gejala pada siklus permainan, penyebab, dan cara memeriksanya}
\label{tab:gejala-siklus}
{\setlength{\tabcolsep}{4pt}\footnotesize%
\begin{tabularx}{\textwidth}{@{}L{3.8cm} L{4.6cm} Y@{}}
\toprule
\headrow \textbf{Gejala} & \textbf{Kemungkinan penyebab} & \textbf{Cara memeriksa} \\
\midrule
Jendela tidak menanggapi dan judulnya tidak berubah &
\code{pygame.event.get()} tidak dipanggil setiap bingkai, atau siklus berhenti
karena kesalahan yang tidak terlihat. &
Periksa pesan kesalahan pada terminal, lalu pastikan pemanggilan
\code{event.get()} berada di dalam perulangan. \\
\addlinespace
Objek meninggalkan jejak di layar &
Layar tidak diisi ulang pada awal tahap penayangan. &
Pastikan \code{screen.fill()} dijalankan sebelum menggambar apa pun pada
setiap bingkai. \\
\addlinespace
Penghitung bingkai selalu bernilai satu &
Variabel keadaan dibuat di dalam perulangan sehingga nilainya kembali ke awal
setiap bingkai. &
Pindahkan pembuatan variabel keadaan ke luar perulangan. \\
\addlinespace
Judul jendela menampilkan laju yang jauh di bawah batas &
Pekerjaan satu bingkai terlalu berat untuk jatah waktunya. &
Ukur dengan program pengukuran, lalu kurangi pekerjaan pada tahap yang paling
banyak memakan waktu. \\
\addlinespace
Prosesor bekerja penuh terus-menerus &
\code{clock.tick()} dipanggil tanpa argumen batas laju. &
Berikan batas laju, misalnya \code{clock.tick(60)}. \\
\bottomrule
\end{tabularx}}
\end{table}

Keempat gejala pertama tidak menimbulkan pesan kesalahan.
Program tetap berjalan, dan tidak ada satu pun aturan bahasa yang dilanggar.
Karena itu, kemampuan membaca gejala menjadi penting: kesalahan seperti ini
hanya dapat ditemukan dengan mengamati perilaku program, bukan dengan
menunggu pesan galat.
